\section{Glass-Box Reasoning Studio: Visualizing Graph Attention and Latent Multi-Hop Inference}

\textbf{Authors}: Bryan Alexander Buchorn  
\textbf{Affili\-ations}: Independent Researcher, Las Vegas, NV, USA  
\textbf{Status}: v1.2.0 (Hardened Enterprise)  
\textbf{Date}: March 2026

---

\subsubsection{Abstract}
The "Black-Box" nature of modern Graph Neural Networks (GNNs) and Transformer-based reasoning systems limits their utility in domains requiring high audit\-ability. We present the \textbf{Glass-Box Reasoning Studio}, an interactive visua\-lization framework designed for the forensic audit of multi-hop Knowledge Graph inference. The Studio reifies the "Reasoning Beam" as a dynamic topological trace, where edges are scaled by their \textbf{Community-Structured Attention (CSA)} weights and nodes are color-coded by their \textbf{DSCF/TSC} community partitions. We introduce a "Forensic Score Breakdown" interface that exposes the latent mathe\-matical signals (semantic similarity, community guidance, and structural centrality) driving each traversal hop, building on found\-ational Explainable AI (XAI) principles \cite{samek2017explainable, ribeiro2016lime, lundberg2017shap}. Furthermore, we describe a real-time "Live Feed" visua\-lization for streaming graphs that animates \textbf{STDP spike events} and the mater\-ialization of speculative causal links. The v1.2.0 release adds adaptive node clustering to support visual scaling for graphs exceeding $10^5$ nodes. Our results show that this interactive "Glass-Box" approach signi\-ficantly reduces the time required for human experts to verify complex AI-generated hypotheses.

\subsubsection{1. Introd\-uction}
Explai\-nability in AI (XAI) has tradi\-tionally focused on post-hoc inter\-pretations of neural weights (e.g., saliency maps). In graph-based reasoning, however, the explanation is the path itself. The Glass-Box Reasoning Studio provides the first integrated environment for visualizing graph attention as a physical, navigatable flow.

\subsubsection{2. Forensic Visual\-ization Methodology}

\paragraph{2.1 The Reasoning Trace}
The Studio implements a path-centric rendering algorithm that isolates the sub-graph involved in a specific query. The attention weight $a(u,v,k)$ is mapped to edge thickness and opacity, allowing the user to visually perceive the "narrowing of the beam" as the AI focuses on likely answers.

\paragraph{2.2 Modal Animations}
For temporal and streaming data, the Studio utilizes high-frequency state updates:
\begin{itemize}
\item \textbf{Potent\-iation}: Edges being stren\-gthened by LTP (SPEC\_003) increase in saturation.
\item \textbf{Drift}: Community boundaries shift smoothly using force-directed layouts to reflect modularity updates (SPEC\_007).

\end{itemize}
\subsubsection{3. Interactive Debugging (v1.2.0)}
The Studio provides a "Dialectical reasoning" mode where users can manually adjust CSA parameters ($\alpha, \beta, \gamma$) via sliders and observe the immediate physical shift in the reasoning beam, providing a "Human-in-the-Loop" (HITL) interface for hyper\-parameter tuning. In v1.2.0, this includes real-time feedback submission to the \textbf{MetaParameterLearner}.

\subsubsection{4. Conclusion}
The Glass-Box Reasoning Studio transforms graph attention from an abstract mathe\-matical construct into a tangible, auditable artifact. By bridging the gap between latent semantic operations and human-readable topologies, it enables the deployment of autonomous reasoning systems in high-stakes envir\-onments.

---
\textbf{References}
\begin{enumerate}
\item Ribeiro, M. T., et al. (2016). "Why Should I Trust You?": Explaining the Predictions of Any Classifier. KDD.
\item Bastian, M., et al. (2009). Gephi: An Open Source Software for Exploring and Manipu\-lating Networks. ICWSM.
\item Hohman, F., et al. (2018). Visual Analytics in Deep Learning: An Interr\-ogative Survey for the Next Frontier. IEEE TVCG.
\item Miller, T. (2019). Explanation in artificial intel\-ligence: Insights from the social sciences. Artificial Intell\-igence.
\item Samek, W., et al. (2017). Explainable AI: Interp\-reting, Explaining and Visualizing Deep Learning. Springer.
\item Lundberg, S. M., \& Lee, S. I. (2017). A Unified Approach to Interp\-reting Model Predictions. NIPS.
\item Buchorn, B. A., \& Sonnet, C. (2026). Interactive Graph Attention in CEREBRUM. SPEC\_012.md.

\end{enumerate}